\begin{afterword}
\summaryhead

The post retells Lewis Carroll's ``What the Tortoise Said to Achilles.'' The Tortoise accepts two premises from Euclid but not the conclusion, so Achilles asks it to accept the rule that the conclusion follows, as a further premise. The Tortoise accepts it and still refuses the conclusion, and each new rule written down calls for another. Douglas Hofstadter's version shows each rule of inference turned into one more string.

The author's lesson is that a mind needs the dynamic of actually drawing the conclusion: a process, not another belief. A rock lacks it, and no argument can supply it; a mind must be ``created already in motion.'' The same holds for action. A mind that believes that pulling a toddler off the tracks is ``fuzzle'' will do nothing unless it implements a dynamic that sends fuzzle plans to action, and beliefs about that dynamic, or proofs that minds with it are more fuzzle, move only a mind that already has the dynamic of adopting better code. So you cannot argue fuzzleness into a rock.

\responsehead

The post makes one point and makes it well. A rule of inference, or a rule for turning a verdict into action, cannot be given to a mind as one more belief, because each new belief needs the rule before it has any effect. The post credits Carroll for the regress and Hofstadter for the notation, and quotes both accurately.

What it adds to Carroll is a reading and an extension, and both have precedents. The reading, that the Tortoise lacks a process no premise can supply, was Gilbert Ryle's in 1946: knowing a rule of inference ``is not possessing a bit of extra information'' but being able to do something. The extension to action, that a belief that something is fuzzle moves nothing without a dynamic that turns fuzzle plans into acts, is Hume's view that reason alone cannot motivate the will. Nick Bostrom later used the same Humean view in support of the orthogonality thesis, that intelligence and final goals can vary independently, and Simon Blackburn had retold Carroll's dialogue about action in 1995 (a paper I could read only in summaries). The post does not claim these points as new, and they support it. The author had also made the modus ponens point before, in ``A Priori.''

In short: a correct point, credited to Carroll and clearly made, whose reading of the dialogue and whose extension to action had been given before by Ryle and by Hume.
\end{afterword}
